%------------------------------------------------------------------------------%
\begin{question}
  Considerando a equação
  \[
    3x^2 + 4xy + 3y^2 - 15 = 0
  \]
  \vspace{-\baselineskip}
  \begin{enumerate}
    \item escreva a equação na forma matricial
    \item diagonalize a matriz da parte quadrática
    \item faça a mudança de variáveis
    \item identifique a cônica
  \end{enumerate}
\end{question}

%------------------------------------------------------------------------------%
\begin{resolution}{Equação Matricial}
  Equação
  \[
    3x^2 + 4xy + 3y^2 - 15 = 0
  \]
  Equação na forma matricial
  \[
    v^tAv + B^tv + f = 0
  \]
  \[
    v = \begin{bmatrix} x \\ y \end{bmatrix}
    \qquad
    A = \begin{bmatrix}
      3 & 2 \\
      2 & 3
    \end{bmatrix}
    \qquad
    B = \begin{bmatrix} 0 \\ 0 \end{bmatrix}
    \qquad
    f = -15
  \]
\end{resolution}

%------------------------------------------------------------------------------%
\begin{resolution}{Diagonalizando $A$ -- Polinômio Característico}
  Polinômio característico
  \begin{align*}
    p(\lambda)
    & = \det(A-\lambda I)
    \\
    & = \begin{vmatrix}
      3-\lambda & 2 \\
      2 & 3-\lambda
    \end{vmatrix}
    \\
    & = (3-\lambda)^2 - 4
    \\
    & = 9 - 6\lambda + \lambda^2 - 4
    \\
    & = \lambda^2 - 6\lambda + 5
  \end{align*}
\end{resolution}

%------------------------------------------------------------------------------%
\begin{resolution}{Diagonalizando $A$ -- Autovalores}
  Equação característica
  \[
    \lambda^2 - 6\lambda + 5 = 0
  \]
  \[
    (\lambda-5)(\lambda-1)=0
  \]
  Autovalores
  \[
    \lambda_1=5
    \qquad
    \lambda_2=1
  \]
\end{resolution}

%------------------------------------------------------------------------------%
\begin{resolution}{Autovetores Associados a $\lambda_1=5$}
  \begin{align*}
    (A-5I)v_1 & = \vec{0} \\
    \begin{bmatrix}
      -2 & 2 \\
      2 & -2
    \end{bmatrix}
    \vetor{x}{y}
    & = \vetor{0}{0}
  \end{align*}
  Portanto
  \[
    -2x+2y=0
    \qquad
    y=x
  \]
  Escolhendo \(x=1\)
  \[
    v_1=\vetor{1}{1}
  \]
\end{resolution}

%------------------------------------------------------------------------------%
\begin{resolution}{Autovetores Associados a $\lambda_2=1$}
  \begin{align*}
    (A-I)v_2 & = \vec{0} \\
    \begin{bmatrix}
      2 & 2 \\
      2 & 2
    \end{bmatrix}
    \vetor{x}{y}
    & = \vetor{0}{0}
  \end{align*}
  Portanto
  \[
    2x+2y=0
    \qquad
    y=-x
  \]
  Escolhendo \(x=1\)
  \[
    v_2=\vetor{1}{-1}
  \]
\end{resolution}

%------------------------------------------------------------------------------%
\begin{resolution}{Matriz Ortogonal}
  Normalizando os autovetores
  \[
    u_1=\frac{1}{\sqrt{2}}\vetor{1}{1}
    \qquad
    u_2=\frac{1}{\sqrt{2}}\vetor{1}{-1}
  \]
  Matriz ortogonal
  \[
    P=\left[u_1\;u_2\right]
    =\frac{1}{\sqrt{2}}
    \begin{bmatrix}
      1 & 1 \\
      1 & -1
    \end{bmatrix}
  \]
  Matriz diagonal
  \[
    D=P^tAP
    =\begin{bmatrix}
      5 & 0 \\
      0 & 1
    \end{bmatrix}
  \]
\end{resolution}

%------------------------------------------------------------------------------%
\begin{resolution}{Mudança de Variáveis}
  Mudança de variáveis
  \[
    v=Pw
    \qquad
    w=\vetor{\mu}{\nu}
  \]
  \[
    \begin{bmatrix}
      x \\ y
    \end{bmatrix}
    =
    \frac{1}{\sqrt{2}}
    \begin{bmatrix}
      1 & 1 \\
      1 & -1
    \end{bmatrix}
    \begin{bmatrix}
      \mu \\ \nu
    \end{bmatrix}
  \]
  Portanto
  \[
    x=\frac{\mu+\nu}{\sqrt{2}}
    \qquad
    y=\frac{\mu-\nu}{\sqrt{2}}
  \]
\end{resolution}

%------------------------------------------------------------------------------%
\begin{resolution}{Equação Transformada}
  Como \(B=\vec{0}\), temos \(E=P^tB=\vec{0}\).

  \begin{align*}
    w^tDw+E^tw+f & = 0
    \\
    \begin{bmatrix}\mu&\nu\end{bmatrix}
    \begin{bmatrix}
      5&0\\
      0&1
    \end{bmatrix}
    \begin{bmatrix}\mu\\\nu\end{bmatrix}
    -15 & = 0
    \\
    5\mu^2+\nu^2 & = 15
  \end{align*}
\end{resolution}

%------------------------------------------------------------------------------%
\begin{resolution}{Identificação}
  Dividindo a equação por \(15\)
  \[
    \frac{\mu^2}{3}+\frac{\nu^2}{15}=1
  \]
  A cônica é uma elipse com centro na origem.

  Semieixos
  \[
    a=\sqrt{15}
    \qquad
    b=\sqrt{3}
  \]

  O eixo maior está na direção de \(\nu\)
\end{resolution}

%------------------------------------------------------------------------------%
\begin{resolution}{Esboço}
  \centering
  \includegraphics{I-08-Rotacao-ex-1}
\end{resolution}
%------------------------------------------------------------------------------%

